We introduce a finite registered geometry built from the 4-regular
arc-transitive graph \(\operatorname{AT4val}[60,6]\), independently
identified as House of Graphs 52002. The construction admits two
complementary presentations. In the forward presentation, a spiral
develops outward from a distinguished centre through successive
registered stages whose initial edge lengths are \(1,4,8\). In the
inverse presentation, the same stages appear as nested three-sided cells
equipped with left-right traversal words. Fifteen registered placements
of the sixty-state carrier form a 900-state closure. The central problem
is to derive the equivalence of the forward and inverse presentations
and determine the orientation group generated by their native
transitions. We investigate whether one complete circuit closes the
projected geometry while inducing a nontrivial central transformation in
its signed lift, so that the lifted geometry closes only after a second
circuit. The target identities
\[
U(2\pi)=-I,\qquad U(4\pi)=I
\]
are therefore treated as geometric holonomy statements, not imported
quantum postulates. The result, if established, gives an exact finite
construction of double rotation from incidence, orientation, and closure
alone. Its connection with half-integer quantum spin is a consequence of
the geometry rather than an assumption used to build it.
